\documentclass[oneside,a4paper,article]{memoir}
\semiisopage
\usepackage{livros_crp}
\usepackage{longtable}
\usepackage{booktabs}
\usepackage{array}
\usepackage{makecell}
\renewcommand\theadalign{cc}
\renewcommand\theadfont{\bfseries}
\renewcommand{\cellalign}{cc}
\renewcommand\theadgape{\Gape[4pt]}
\renewcommand\cellgape{\Gape[4pt]}
\title{Estudos técnicos preliminares para contratação de agência de comunicação}
\author{Coordenação de Comunicação do CRP~SP}
\date{\today}
\begin{document}
\pagestyle{crp}
\maketitle

O presente documento caracteriza a primeira etapa da fase de
planejamento e apresenta os devidos estudos para a contratação de
solução que atenderá à necessidade abaixo especificada.

O objetivo principal é estudar detalhadamente a necessidade e
identificar no mercado a melhor solução para supri-la, em observância às
normas vigentes e aos princípios que regem a Administração Pública.
\vspace{3\baselineskip}

\chapter{Descrição da necessidade}

\section{Competências da unidade de comunicação}

Em conformidade com suas atribuições fundamentais estabelecidas na Lei
Federal nº~5.766/1971, que incluem orientar, fiscalizar, disciplinar,
normatizar e proteger o exercício profissional de psicólogos e
psicólogas, o Conselho Regional de Psicologia da 6ª Região --- CRP~SP,
atende às demandas provenientes de suas comissões, áreas administrativas
e territórios. A equipe de comunicação, conforme estipulado no artigo da
Resolução CRP N. 003/2022, que regula a Estrutura Administrativa e o
Quadro de Gestão do CRP~SP, possui entre suas competências:

\begin{enumerate}
	\item
	      planejar, organizar, coordenar e divulgar as ações referentes à
	      integração das estratégias institucionais com a categoria e a
	      sociedade;
	\item
	      construir, desenvolver e implantar estratégias com foco na aproximação
	      com a categoria e a sociedade;
	\item
	      garantir a divulgação das atividades e posicionamentos institucionais
	      do CRP-06 para a categoria e a sociedade.
\end{enumerate}

Regularmente, a Coordenação de Comunicação elabora e organiza campanhas
e eventos direcionados à categoria, que servem como veículos para
difundir orientações, regulamentos e informações do Sistema Conselhos de
Psicologia para as psicólogas e para a sociedade brasileira. O objetivo
é promover o engajamento e o incentivo em relação às temáticas, ações e
mensagens da Psicologia, embasadas cientificamente e comprometidas
eticamente, contribuindo assim para consolidar o CRP~SP como uma
referência na esfera pública.

Para atingir os objetivos da Resolução CRP nº~003/2022, é fundamental
adotar uma estratégia de comunicação eficaz, assegurando a transparência
das ações da autarquia e fortalecendo sua imagem perante psicólogas e a
sociedade, além de aumentar sua capacidade de disseminação de
informações e ações. A comunicação digital poderá contribuir para o
alcance desses objetivos, uma vez que as estratégias de comunicação
tradicionais estão cada vez mais utilizando plataformas e tecnologias
digitais para atender os objetivos de comunicação, tornando a estratégia
mais completa e efetiva.

\section{Alcance da categoria e difusão da Psicologia por meio	dos canais do Conselho e de veículos da	imprensa}

Em 2024, o CRP~SP ultrapassou a marca de 200 mil profissionais inscritos
(sendo quase 150 mil registros ativos) em 11 territórios, buscando a
divulgação e promoção de conteúdos relevantes e informativos acerca da
atuação ética e legal da categoria.

Considerando a importância estratégica das divulgações de ações em
diversos canais de comunicação e o diálogo constante do CRP~SP com a
categoria, entidades, poder público e sociedade, bem como a crescente
demanda por produção e disseminação de conteúdo associada às atividades
do Conselho, torna-se imprescindível a contratação de uma empresa
especializada em comunicação digital. Essa contratação permitirá uma
abordagem mais abrangente e eficaz na promoção das atividades e na
construção de uma imagem sólida e representativa do CRP~SP, além de
garantir uma interação mais dinâmica e eficiente com os diversos
públicos envolvidos. A comunicação digital oferece recursos e
ferramentas inovadoras que possibilitam uma maior visibilidade, alcance
e impacto das ações do Conselho, contribuindo assim para o
fortalecimento de sua atuação e para a consolidação de sua posição como
referência no campo da Psicologia e no cenário público em geral.

No que tange ao item 2.1 do Processo de Planejamento operacional do CRP
2024, dos Resultados Esperados, a consolidação do CRP~SP como
instituição de referência na luta contra violências estruturais, em
colaboração com os movimentos sociais, para a construção democrática de
políticas públicas, demanda uma comunicação eficaz. Uma empresa de
comunicação digital pode contribuir para promover campanhas de
conscientização, divulgar informações relevantes e engajar a sociedade
em ações de combate às violências estruturais, por meio de estratégias
\emph{on-line} que alcancem um público mais amplo e diversificado.

A utilização da comunicação digital é essencial para alcançar os
objetivos da comunicação institucional, especialmente em São Paulo, o
segundo estado brasileiro com maior acesso à internet em domicílios,
conforme dados do Instituto Brasileiro de Geografia e Estatística
(IBGE), além de ser um dos líderes em acessos móveis no país.

No contexto do CRP~SP, as plataformas de comunicação digital têm um
alcance significativo: em 2023, mais de 4.000.000 (quatro milhões) de
usuários acessaram o portal do órgão na internet. Nas redes sociais,
como Facebook e Instagram, o CRP~SP conta com uma audiência expressiva,
com aproximadamente 158.000 (cento e cinquenta e oito mil) seguidores no
Facebook e mais de 100.000 (cem mil) no Instagram. As publicações
alcançaram mais de 575.000 (quinhentas e setenta e cinco mil) pessoas no
Facebook e mais de 3.000.000 (três milhões) no Instagram.

Esses resultados são fruto do planejamento das ações de comunicação
digital, da produção de conteúdo relevante e da padronização da
identidade visual, o que torna as informações mais atrativas e promove
maior engajamento do público-alvo. Nesse cenário, um serviço
especializado se mostra fundamental, dada a constante evolução e
surgimento de novas plataformas no universo da comunicação digital, além
das mudanças nos algoritmos que impactam na distribuição do conteúdo.

\section{Previsão no plano de contratações anual}

O Conselho Regional de Psicologia da 6ª Região --- CRP~SP, por ser uma
autarquia federal de natureza \emph{sui generis}, não está obrigado à
elaboração de Plano Anual de Contratações.

\chapter{Requisitos da contratação}

Os produtos e serviços constantes deste ETP serão executados e entregues
continuamente, \textbf{mediante demanda, na forma de execução indireta,
	sob o regime de empreitada por preço unitário}.

A empresa será contratada por meio de \textbf{concorrência, do tipo
	melhor técnica}, sob a égide da Lei nº~8.666/1993 e da Instrução
Normativa SEGES/MPDG nº~05/2017.

Subsidiariamente, devem ser observadas as regras estabelecidas na Lei
nº~12.232/2010, aplicáveis a este objeto, nos termos do Acórdão TCU
nº~6.227/2016-2ª Câmara e na Instrução Normativa Secom relativa à
matéria.

%\centering
	\begin{table}
	\caption{Estimativa de quantidades.\label{tab:estimativa}}%\footnotesize
%^	\footnotesize
\begin{center}
	\begin{tabular}{c c c c}
		\toprule\noalign{}
		\thead{Item} & \thead{Descrição/\\especificação} & \thead{Qtd. total\\estimada\\(12 meses)} & \thead{Unidade}   \\
		\midrule\noalign{}
		\makecell{\textbf{1}}    & \makecell{Criação de identidade visual}      & 1 & \makecell{Unidade/\\manual} \\
		\makecell{\textbf{2}}    & \makecell{Adaptação ou replicação\\de postagem (modelos)}   & \makecell{12} & \makecell{Mensal} \\
		\makecell{\textbf{3}}    & \makecell{Planejamento estratégico de\\ comunicação digital} & \makecell{1} & \makecell{Unidade/\\relatório} \\
		\makecell{\textbf{4}}    & \makecell{Análise e matriz estratégica }     & \makecell{1} & \makecell{Unidade/\\relatório}\\
		\makecell{\textbf{5}}    & \makecell{Diagnóstico de conteúdo} & \makecell{1} & \makecell{Unidade/\\relatório} \\
		\makecell{\textbf{6}}    & \makecell{Treinamento de equipe} & \makecell{1} & \makecell{Treinamento}\\
		\makecell{\textbf{7}}    & \makecell{Planejamento de conteúdo} & \makecell{1} & \makecell{Unidade/\\relatório} \\
		\makecell{\textbf{8} }   & \makecell{Produção de conteúdo} & \makecell{12} & \makecell{Mensal}\\
		\makecell{\textbf{9}} & \makecell{Produção de vídeo\\(duração de até 2 minutos)}    & \makecell{12} & \makecell{Mensal}\\
		\makecell{\textbf{10}}   & \makecell{Produção de animação\\ (duração de até 2 minutos)} & \makecell{12 } & \makecell{Mensal } \\
		\makecell{\textbf{11}}   & \makecell{Monitoramento e emissão\\ de relatórios mensais}   & \makecell{12} & \makecell{Mensal} \\
		\makecell{\textbf{12}}   & \makecell{Relatório de\\ desempenho anual  }   & \makecell{1} & \makecell{Unidade/\\relatório}\\
		\makecell{\textbf{13} }  & \makecell{Impulsionamento e\\ gestão de tráfego } & \makecell{12} & \makecell{Mensal}  \\
		\bottomrule\noalign{}
	\end{tabular}
\end{center}
\end{table}

\chapter{Levantamento de mercado}

Do levantamento realizado no mercado, identifica-se a
existência de duas soluções possíveis (Art. 7°, inciso III da IN
40/2020):

\begin{itemize}
	\item \textbf{Solução 1:} Realizar internamente as demandas descritas no item 1;
	\item \textbf{Solução 2:} Contratar de empresa de comunicação digital.
\end{itemize}

A \textbf{solução 1} citada é manifestamente inviável, pois o Conselho
Regional de Psicologia da 6ª Região --- CRP~SP não detém de condições
técnicas e de pessoal adequadas para a execução dos trabalhos.

A equipe de comunicação do CRP~SP é responsável por gerenciar o conteúdo
de todas as unidades do órgão, concentrando informações de interesse
público e disseminando-as de forma eficiente. Portanto é inviável
agregar mais uma demanda, pois prejudicaria o desempenho de suas
competências na área de comunicação de maneira eficaz e abrangente.

Profissionais que atuam na comunicação digital são especialistas
responsáveis por planejar, desenvolver e executar estratégias de
comunicação utilizando meios digitais e ferramentas específicas para
alcançar determinados objetivos organizacionais. Suas funções e áreas de
atuação podem variar dependendo das necessidades e dos objetivos
específicos de cada organização, mas geralmente incluem:

\begin{itemize}
	\item
	      \textbf{estratégia digital:} desenvolvimento de planos estratégicos de
	      comunicação digital alinhados aos objetivos da organização,
	      identificando o público-alvo, definindo canais de comunicação e
	      estabelecendo métricas de sucesso.
	\item
	      \textbf{produção de conteúdo:} criação de conteúdo relevante e atrativo para
	      ser veiculado nos canais digitais da organização, como redes sociais,
	      \textit{blogs}, \textit{sites}, \textit{e-mail marketing}, vídeos, \textit{podcasts}, entre outros.
	\item
	      \textbf{gestão de redes sociais:} administração e monitoramento das redes
	      sociais da organização, incluindo o planejamento de postagens,
	      interação com os seguidores, análise de métricas e elaboração de
	      relatórios de desempenho.
	\item
	      \textbf{SEO (\textit{search engine optimization}):} otimização de conteúdo para
	      mecanismos de busca com o objetivo de aumentar a visibilidade e o
	      tráfego orgânico nos \textit{sites} e páginas da organização.
	\item
	      \textbf{gestão de mídias pagas:} criação e gerenciamento de campanhas de
	      publicidade  \textit{on-line}, incluindo anúncios em redes sociais, Google Ads,
	      entre outros, visando aumentar o alcance e a visibilidade da marca.
	\item
	      \textbf{análise de dados:} Coleta e análise de dados sobre o desempenho das
	      estratégias de comunicação digital, utilizando ferramentas de análise
	      para avaliar o impacto das ações e identificar oportunidades de
	      melhoria.
\end{itemize}

Um dos principais desafios enfrentados pela equipe interna é a falta de
recursos especializados e ferramentas específicas para atender às
demandas complexas da comunicação digital. Embora sejam profissionais
capacitados, a equipe pode enfrentar limitações técnicas e de expertise
na implementação de estratégias avançadas de comunicação digital, o que
pode comprometer a eficácia e a eficiência das ações.

Além disso, a sobrecarga de trabalho da equipe interna, que já é
responsável por diversas outras atividades relacionadas à comunicação do
CRP~SP, pode dificultar a execução adequada das demandas de comunicação
digital. A falta de tempo e de foco exclusivo nessas atividades pode
resultar em atrasos na entrega de conteúdo, falta de planejamento
estratégico e, consequentemente, na não obtenção dos resultados
esperados.

A \textbf{solução 2}, contratação de uma empresa de comunicação digital,
atingiria os eixos dos resultados esperados da Gestão do XVII Plenário
do CRP~SP, especialmente o Eixo 2 (processo de planejamento operacional
do CRP 2024), que diz respeito ao diálogo da Psicologia com a sociedade
brasileira e suas relações com a democracia e direitos humanos.

A contratação de uma empresa especializada em comunicação digital se
justifica pelo fato de que no quadro de trabalhadores do CRP~SP não
existem profissionais com as competências e habilidades necessárias para
desempenhar essas funções de forma eficaz. Uma empresa especializada
traz consigo uma equipe multidisciplinar com experiência em diversas
áreas da comunicação digital, garantindo uma abordagem abrangente e
alinhada com as melhores práticas do mercado.

Além disso, uma empresa especializada possui recursos e ferramentas
específicas para realizar análises de mercado, monitorar o desempenho
das estratégias adotadas e implementar soluções inovadoras, contribuindo
assim para o alcance dos objetivos de comunicação do CRP~SP, incluindo a
promoção do diálogo com a sociedade, o fortalecimento da marca e a
disseminação de informações relevantes para a categoria profissional.

%Analisando as questões que impedem a realização através do item 1, sob os aspectos da economia e eficiência, consideramos que a solução 2 (contratação de uma empresa de comunicação digital) é a mais indicada para atender às necessidades do Conselho Regional de Psicologia da 6ª Região --- CRP~SP, uma vez que essas empresas são especializadas na área.
%
Considerando esses aspectos, a solução mais indicada para contemplar as
necessidades do CRP~SP é a contratação de uma empresa especializada em
comunicação digital. As empresas desse segmento são compostas por
profissionais altamente qualificados e experientes, com expertise em
diversas áreas da comunicação digital, como estratégia, produção de
conteúdo, gestão de redes sociais, SEO, entre outras.

É relevante ressaltar que a comunicação digital é uma das formas mais
econômicas e sustentáveis de comunicação, oferecendo uma relação
custo-benefício mais vantajosa do que muitas mídias tradicionais e
offline. Além disso, as estratégias digitais permitem a segmentação do
público-alvo, garantindo um maior retorno sobre o investimento e uma
melhor mensuração dos resultados obtidos.

No que se refere ao item 2.2, que aborda a promoção de ações que
engajem a categoria na defesa dos Direitos Humanos e do Estado
democrático de direito, a comunicação digital desempenha um papel
crucial. Uma empresa especializada pode apoiar na criação de conteúdos
que sensibilizem e mobilizem os psicólogos e psicólogas para a
importância da atuação profissional em consonância com os princípios dos
Direitos Humanos e da democracia, promovendo debates, eventos e
iniciativas  \textit{on-line} que estimulem a reflexão e a participação ativa da
categoria. Assim, a contratação de uma empresa de comunicação digital se
alinha diretamente aos objetivos estabelecidos pelo XVII Plenário do
CRP~SP, pois contribui para fortalecer o diálogo da Psicologia com a
sociedade brasileira, ampliando o alcance das ações e promovendo o
engajamento da categoria na defesa dos Direitos Humanos e da democracia.

A contratação da empresa de comunicação digital também se justifica
à luz do Eixo 3 (``do exercício profissional''), mais especificamente no
item 3.4. ``dos resultados esperados da gestão do XVII Plenário do CRP~SP'',
destaca a importância de executar projetos que contemplem a orientação
do exercício ético da Psicologia e monitorar as diferentes práticas da
profissão, promovendo o diálogo sobre condutas emancipatórias. Nesse
sentido, uma empresa de comunicação digital pode desempenhar um papel
fundamental ao desenvolver conteúdos e campanhas que orientem os
profissionais sobre condutas éticas, promovam a reflexão sobre práticas
inclusivas e emancipatórias, e incentivem a adesão a princípios éticos
no exercício da Psicologia.

Por meio de estratégias digitais, como a criação de materiais
educativos e \textit{webinars} e a divulgação de casos e exemplos de boas
práticas profissionais, a empresa de comunicação digital pode contribuir
para a disseminação de informações relevantes e atualizadas sobre ética
na Psicologia, favorecendo o aprimoramento contínuo dos profissionais e
o fortalecimento da categoria como um todo.

Dessa forma, a contratação de uma empresa de comunicação digital se
mostra essencial para atender às demandas do Eixo 3 --- Do Exercício
Profissional, ao proporcionar um canal eficaz de comunicação que promova
o diálogo, a reflexão e a orientação ética no contexto da prática
psicológica.

%\begin{landscape}
\chapter{Estimativa do preço da contratação}

Estima-se que o valor seria \textbf{R\$ 1.291.000,00}\ldots isso é uma média de valores que obtivemos através de consulta de orçamento etc \ldots{}

Fizemos cotação de maneira formal com três empresas, conforme tabelas \ref{llyc} a \ref{fb}.

%\begin{landscape}
	\begin{table}
		\caption{Razão social: LLYC --- Llorente \& Cuenca do Brasil Consultores de Comunicação Ltda. CNPJ: 09.429.171/0001-05.\label{llyc}}
\begin{center}\footnotesize
	\begin{tabular}{c c c c c c}
		\toprule\noalign{}
		\thead{Item} & \thead{Descrição/\\especificação} & \thead{Qtd. total\\estimada\\(12 meses)} & \thead{Unidade} & \thead{Valor\\unitário} & \thead{Valor\\total}  \\
		\midrule\noalign{}
			\makecell{\textbf{1} }   & \makecell{Criação de\\ identidade visual}      & \makecell{1} & \makecell{Unidade/\\manual}    & \makecell{R\$ 25.500,00} & \makecell{R\$ 25.500,00} \\
			\makecell{\textbf{2}}   & \makecell{Adaptação ou replicação\\de postagem (modelos)}   & \makecell{12} & \makecell{Mensal} & \makecell{R\$ 6.700,00} & \makecell{R\$ 80.400,00} \\
			\makecell{\textbf{3}}    & \makecell{Planejamento estratégico\\ de comunicação digital} & 1 & \makecell{Unidade/\\relatório} & \makecell{R\$ 28.500,00} & \makecell{R\$ 28.500,00} \\
			\makecell{\textbf{4}}    & \makecell{Análise e matriz estratégica}      & \makecell{1 } & \makecell{Unidade/\\relatório} & \makecell{R\$ 27.530,00} & \makecell{R\$ 27.530,00} \\
			\makecell{\textbf{5}}    & \makecell{Diagnóstico de conteúdo} & \makecell{1} & \makecell{Unidade/\\relatório} & \makecell{R\$ 31.850,00} & \makecell{R\$ 31.850,00 } \\
			\makecell{\textbf{6}}    & \makecell{Treinamento de equipe} & \makecell{1} & \makecell{Treinamento} & \makecell{R\$ 25.000,00} & \makecell{R\$ 25.000,00} \\
			\makecell{\textbf{7}}    & \makecell{Planejamento de conteúdo} & \makecell{1} & \makecell{Unidade/\\relatório} & \makecell{R\$ 29.128,00} &\makecell{ R\$ 29.128,00} \\
			\makecell{\textbf{8}}    & \makecell{Produção de conteúdo} & \makecell{12} & \makecell{Mensal} & \makecell{R\$ 13.560,00} & R\$ 162.720,00 \\
			\makecell{\textbf{9}}   & \makecell{Produção de vídeo\\(duração de até 2 minutos)}    & \makecell{12} & \makecell{Mensal} & \makecell{R\$ 12.250,00} & \makecell{R\$ 147.000,00 } \\
			\makecell{\textbf{10}}   & \makecell{Produção de animação\\ (duração de até 2 minutos)} & \makecell{12} & \makecell{Mensal} & \makecell{R\$ 15.420,00} & \makecell{R\$ 185.040,00} \\
			\makecell{\textbf{11}}   & \makecell{Monitoramento e emissão\\de relatórios mensais}   & \makecell{12} & \makecell{Mensal} & \makecell{R\$ 5.230,00} & \makecell{R\$ 62.760,00} \\
			\makecell{\textbf{12}}   & \makecell{Relatório de desempenho anual}    & \makecell{1} & \makecell{Unidade/\\relatório} & \makecell{R\$ 12.972,00} & \makecell{R\$ 12.972,00} \\
			\makecell{\textbf{13}}   & \makecell{Impulsionamento e gestão\\de tráfego} & \makecell{12} & \makecell{Mensal} & \makecell{R\$ 5.000,00} & \makecell{R\$ 60.000,00 } \\
						\midrule\noalign{}
			& & & &\makecell{\textbf{Valor total}} &\makecell{\textbf{R\$ 790.560,00}}\\
\bottomrule\noalign{}
		\end{tabular}
	\end{center}
	\end{table}
\vfill\clearpage
%\end{landscape}


	\begin{table}\footnotesize
		\caption{Razão social: JG Comunicações Brasil Ltda. --- Jeffrey Group. CNPJ:~46.504.820/0001-11.\label{jg}}
\begin{center}\footnotesize
	\begin{tabular}{c c c c c c}
		\toprule\noalign{}
		\thead{Item} & \thead{Descrição/\\especificação} & \thead{Qtd. total\\estimada\\(12 meses)} & \thead{Unidade} & \thead{Valor\\unitário} & \thead{Valor\\total}  \\
		\midrule\noalign{}
		\makecell{\textbf{1} }   & \makecell{Criação de\\ identidade visual}      & \makecell{1} & \makecell{Unidade/\\manual}    & \makecell{R\$ 21.657,50 (?)} & \makecell{R\$ } \\
		\makecell{\textbf{2}}   & \makecell{Adaptação ou replicação\\de postagem (modelos)}   & \makecell{12} & \makecell{Mensal} & \makecell{R\$ } & \makecell{R\$ } \\
		\makecell{\textbf{3}}    & \makecell{Planejamento estratégico\\ de comunicação digital} & 1 & \makecell{Unidade/\\relatório} & \makecell{R\$ } & \makecell{R\$ } \\
		\makecell{\textbf{4}}    & \makecell{Análise e matriz estratégica}      & \makecell{1 } & \makecell{Unidade/\\relatório} & \makecell{R\$ } & \makecell{R\$ } \\
		\makecell{\textbf{5}}    & \makecell{Diagnóstico de conteúdo} & \makecell{1} & \makecell{Unidade/\\relatório} & \makecell{R\$ } & \makecell{R\$ } \\
		\makecell{\textbf{6}}    & \makecell{Treinamento de equipe} & \makecell{1} & \makecell{Treinamento} & \makecell{R\$ } & \makecell{R\$ } \\
		\makecell{\textbf{7}}    & \makecell{Planejamento de conteúdo} & \makecell{1} & \makecell{Unidade/\\relatório} & \makecell{R\$ } &\makecell{ R\$ } \\
		\makecell{\textbf{8}}    & \makecell{Produção de conteúdo} & \makecell{12} & \makecell{Mensal} & \makecell{R\$ } & R\$ \\
		\makecell{\textbf{9}}   & \makecell{Produção de vídeo\\(duração de até 2 minutos)}    & \makecell{12} & \makecell{Mensal} & \makecell{R\$ } & \makecell{R\$ } \\
		\makecell{\textbf{10}}   & \makecell{Produção de animação\\ (duração de até 2 minutos)} & \makecell{12} & \makecell{Mensal} & \makecell{R\$ } & \makecell{R\$ } \\
		\makecell{\textbf{11}}   & \makecell{Monitoramento e emissão\\de relatórios mensais}   & \makecell{12} & \makecell{Mensal} & \makecell{R\$ } & \makecell{R\$ } \\
		\makecell{\textbf{12}}   & \makecell{Relatório de desempenho anual}    & \makecell{1} & \makecell{Unidade/\\relatório} & \makecell{R\$ } & \makecell{R\$ } \\
		\makecell{\textbf{13}}   & \makecell{Impulsionamento e gestão\\de tráfego} & \makecell{12} & \makecell{Mensal} & \makecell{R\$ } & \makecell{R\$ } \\
		\midrule\noalign{}
		& & & &\makecell{\textbf{Valor total}} &\makecell{\textbf{R\$ }}\\
		\bottomrule\noalign{}
	\end{tabular}
\end{center}
\end{table}
\vfill\clearpage


	\begin{table}\footnotesize
		\caption{Razão social: Future Brand. CNPJ:\label{fb}}
\begin{center}\footnotesize
	\begin{tabular}{c c c c c c}
		\toprule\noalign{}
		\thead{Item} & \thead{Descrição/\\especificação} & \thead{Qtd. total\\estimada\\(12 meses)} & \thead{Unidade} & \thead{Valor\\unitário} & \thead{Valor\\total}  \\
		\midrule\noalign{}
		\makecell{\textbf{1} }   & \makecell{Criação de\\ identidade visual}      & \makecell{1} & \makecell{Unidade/\\manual}    & \makecell{R\$ 70.000,00} & \makecell{R\$ 70.000,00} \\
		\makecell{\textbf{2}}   & \makecell{Adaptação ou replicação\\de postagem (modelos)}   & \makecell{12} & \makecell{Mensal} & \makecell{R\$ 2.000,00} & \makecell{R\$ 24.000,00} \\
		\makecell{\textbf{3}}    & \makecell{Planejamento estratégico\\ de comunicação digital} & 1 & \makecell{Unidade/\\relatório} & \makecell{R\$ 60.000,00} & \makecell{R\$ 60.000,00} \\
		\makecell{\textbf{4}}    & \makecell{Análise e matriz estratégica}      & \makecell{1 } & \makecell{Unidade/\\relatório} & \makecell{R\$ 204.000,00} & \makecell{R\$ 204.000,00} \\
		\makecell{\textbf{5}}    & \makecell{Diagnóstico de conteúdo} & \makecell{1} & \makecell{Unidade/\\relatório} & \makecell{R\$ 90.000,00} & \makecell{R\$ 90.000,00} \\
		\makecell{\textbf{6}}    & \makecell{Treinamento de equipe} & \makecell{1} & \makecell{Treinamento} & \makecell{R\$ 25.000,00} & \makecell{R\$ 25.000,00} \\
		\makecell{\textbf{7}}    & \makecell{Planejamento de conteúdo} & \makecell{1} & \makecell{Unidade/\\relatório} & \makecell{R\$ 30.000,00} &\makecell{ R\$ 360.000,00} \\
		\makecell{\textbf{8}}    & \makecell{Produção de conteúdo} & \makecell{12} & \makecell{Mensal} & \makecell{R\$ 2.000,00} & \makecell{R\$ 288.000,00} \\
		\makecell{\textbf{9}}   & \makecell{Produção de vídeo\\(duração de até 2 minutos)}    & \makecell{12} & \makecell{Mensal} & \makecell{\scriptsize15\% de gestão\\\scriptsize sobre valor de\\\scriptsize produtora parceira} & \makecell{---} \\
		\makecell{\textbf{10}}   & \makecell{Produção de animação\\ (duração de até 2 minutos)} & \makecell{12} & \makecell{Mensal} & \makecell{R\$ 15.000,00} & \makecell{R\$ 15.000,00} \\
		\makecell{\textbf{11}}   & \makecell{Monitoramento e emissão\\de relatórios mensais\\(mínimo de seis meses\\de contratação}   & \makecell{12} & \makecell{Mensal} & \makecell{R\$ 30.000,00} & \makecell{R\$ 360.000,00} \\
		\makecell{\textbf{12}}   & \makecell{Relatório de desempenho anual}    & \makecell{1} & \makecell{Unidade/\\relatório} & \makecell{R\$ 35.000,00} & \makecell{R\$ 35.000,00} \\
		\makecell{\textbf{13}}   & \makecell{Impulsionamento e gestão\\de tráfego} & \makecell{12} & \makecell{Mensal} & \makecell{R\$ 8.000,00} & \makecell{R\$ 96.000,00} \\
		\midrule\noalign{}
		& & & &\makecell{\textbf{Valor total}} &\makecell{\textbf{R\$ 1.792.000,00}}\\
		\bottomrule\noalign{}
	\end{tabular}
\end{center}
\end{table}
\normalsize

\chapter{Descrição da solução como um todo}

As demandas serão atendidas por meio da combinação dos produtos e
serviços mais adequados para apoiar o CRP~SP na superação de seus
desafios e alcance dos seus objetivos de comunicação, abrangendo:

\begin{enumerate}
	\item \textbf{produtos e serviços essenciais:} contemplam a expertise básica da
	      empresa especializada em comunicação digital na execução do objeto do
	      contrato, sendo os itens previamente especificados e precificados pela
	      Autarquia, com os respectivos quantitativos estimados de execução
	      citados no item 4 deste documento.\newline Os produtos e serviços essenciais contemplam as necessidades
 elementares do CRP~SP, de acordo com as categorias abaixo:
    \begin{enumerate}
	    \item
	    \textit{Design}
 \begin{enumerate}
	 \item
	 Criação de identidade visual

 Desenvolvimento de um conceito criativo aplicável em diversas campanhas de comunicação e produtos digitais. Elaboração de uma identidade visual baseada no logotipo e mote da gestão existentes, acompanhada de um guia abrangente para a aplicação consistente da marca. Inclusão de paleta de cores, seleção tipográfica e demais elementos visuais destinados aos canais digitais, como \textit{site}, redes sociais e \textit{e-mail marketing}, além de materiais impressos. Manual da identidade contendo as diretrizes de utilização dos elementos de \textit{design} para garantir coesão visual.

 Proposta Visual para as Redes Sociais, \textit{site} e \textit{e-mail marketing} do CRP~SP: desenvolvimento de uma proposta visual alinhada à nova identidade visual para as redes sociais.
 \end{enumerate}
      \item
	     Adaptação ou replicação de postagem (modelos)
 \begin{enumerate}
 \item
 Elaboração de modelos editáveis com derivação de \textit{layouts} e alterações necessárias para adequação a diferentes plataformas digitais. Estudo de adequação do estilo da fonte, tamanho e cores, se necessário. Os modelos incluirão \textit{cards} e arquivos para redes sociais e \textit{site}, além de um modelo de apresentações institucionais/relatórios e \textit{e-mail marketing} (boletins,  informativos etc.), proporcionando uma abordagem unificada e impactante.
 \end{enumerate}
    \end{enumerate}
	\item
	\textbf{planejamento estratégico}
		\begin{enumerate}
			\item
	      	Planejamento Estratégico de Comunicação Digital

			Desenvolvimento de um planejamento estratégico para a presença digital conforme o objeto solicitado. A partir da análise e principais pontos do mapeamento, fornecimento de recomendações de ações para cada propriedade digital e gestão da marca, boas práticas, posicionamento e sugestões de iniciativas relacionadas ao órgão/tema. Recomendação de formas de atuação e adequação para melhorias das referidas propriedades digitais. Determinação de campanhas regulares, incluindo site (SEO), e-mail marketing, impulsionamento de posts e iniciativas de conteúdo para manter a marca visível e atrativa.

			\item
	      	Análise e matriz estratégica
				\begin{enumerate}
					\item
	      			Mapeamento de Presença Digital: identificação dos principais assuntos tratados pela pasta, identidade percebida e compartilhada pelo público-alvo, além das necessidades de comunicação. Análise editorial das propriedades digitais, presença em ferramentas de busca e atuação nas redes sociais.
					\item
	      			Acompanhamento da imagem do órgão/tema em veículos  \textit{on-line} e redes sociais, com monitoramento de repercussão, reputação, evolução de sentimento e informações estratégicas para a tomada de decisões. Alertas para situações de possível impacto devem ser providenciados.
					\item
	      			Busca de Melhores Práticas e Desempenho: pesquisa e identificação de melhores práticas e desempenho relacionados à presença digital.
					\item
	      			Definição do Objetivo da Presença Digital: estabelecimento claro do objetivo da presença digital, alinhado aos propósitos do órgão/tema.
					\item
	      			Indicação dos Canais e Propriedades Digitais: identificação dos canais de atuação e propriedades digitais mais adequados para alcançar os objetivos.
					\item
	      			Desenvolvimento da Matriz/Direcionamento Estratégico: consolidação do objetivo da presença digital, delineando estratégias e direcionamentos.
					\item
	      			Definição de indicadores de avaliação de performance estratégica, sujeitos a revisão conforme os objetivos sazonais.
					\item
					Recomendação de parceiros, recursos, gestores e ferramentas essenciais para o sucesso do projeto.
					\item
	      			Benchmark: comparativo estratégico com outros conselhos de classe e órgãos públicos similares ao CRP~SP.
					\item
					Identificação e Estudo dos Públicos: mapeamento e estudo dos públicos, considerando desejos, insumos, potencialidades, fraquezas, oportunidades e ameaças.
	      	\end{enumerate}
		\end{enumerate}
	\item
		\textbf{Diagnóstico de Conteúdo}
			\begin{enumerate}
				\item
	      		Análise detalhada de 4 (quatro) a 6 (seis) canais, considerando sua    eficácia e adequação aos objetivos.
				\item
	      		Mapeamento do Conteúdo: identificação e mapeamento do conteúdo disponibilizado nos canais de comunicação digital do órgão/tema.
				\item
	      		Análise Editorial Aprofundada: avaliação aprofundada do ambiente de comunicação digital do órgão/tema, considerando a qualidade editorial, melhores horários de postagem etc.
				\item
	      		Análise Imagética: análise da comunicação visual utilizada no ambiente digital do órgão/tema.
				\item
	      		\textit{Benchmark}: comparativo estratégico com outros conselhos de classe e órgãos públicos similares ao CRP~SP.
			\end{enumerate}
    \item
 \textbf{Treinamento}

		Treinamento dos membros da equipe de comunicação para assegurar uma comunicação eficaz da nova identidade visual e linguagem. Garantia de que toda a equipe compreenda e aplique as diretrizes de marca estabelecidas.
	\item
	    \textbf{Planejamento Tático}
			\begin{enumerate}
				\item
	      		Planejamento de Conteúdo
				\item
	      		Desenvolvimento de um mapa de conteúdo abrangente, estratégias e campanhas de comunicação direcionadas a diversos públicos relacionados ao tema do órgão. Publicações planejadas em diferentes plataformas.
				\item
	      		Estabelecimento de diretrizes editoriais para criação de títulos, chamadas, tratamento de textos e aplicação de políticas de tagueamento com base no Manual de Linguagem do CRP~SP.
				\item
	      		Criação de um fluxo contínuo de conteúdo para redes sociais, blog, site e outros canais digitais, garantindo uma presença ativa e relevante alinhada com o mote da gestão.
			\end{enumerate}

	\item
	    \textbf{Conteúdo}
			\begin{enumerate}
				\item
	      		Produção de Conteúdo

				Produção e publicação de conteúdo para plataformas digitais (redes sociais, site e e-mail marketing) do CRP~SP. Inclui criação de texto, edição de imagens, tagueamento e publicação de diversos formatos como imagem, texto, vídeos, animações, gifs e infográficos.

				\item
	    		Produção de Vídeo (Duração de até 2 minutos)
				Produção de vídeo conforme briefing e roteiro aprovados, incluindo roteirização, captação de imagens, direção de arte, edição e sonorização, além de adaptação para as diferentes redes. Cessão de direitos autorais e de imagem em arquivo texto. O custo deve abranger todos os processos.
				\item
	    		Produção de Animação (duração de até 2 minutos)
				Produção de vídeo com utilização de técnica de animação, elaborado a partir de briefing e roteiro previamente aprovados. Os profissionais envolvidos devem ceder o uso de direito autoral e de imagem em arquivo texto. O custo deve prever roteirização, direção de arte, edição e sonorização, além de adaptação para as diferentes redes. Os materiais deverão ser submetidos à aprovação em todos os estágios.
			\end{enumerate}
	\item
	    \textbf{Monitoramento}
			\begin{enumerate}
				\item
	      		Monitoramento e Emissão de Relatórios Mensais

				Acompanhamento diário das redes sociais do CRP~SP e páginas da área de atuação. Relatórios mensais incluindo volume total de menções, regionalização e origens de menções, melhores horários de postagens, temas mais comentados, oportunidades para novas ações, rankings de assuntos, saúde dos temas e percepção do público.

				\item
	      		Relatório de Desempenho Anual

				Análise consolidada do desempenho de propriedade digital do CRP~SP, utilizando informações de relatórios de BI, desempenho de redes e monitoramento de redes sociais.
				\item
	      		Gestão de Tráfego e Impulsionamento

				Impulsionamento de publicações em redes sociais e gestão de tráfego para otimização da verba e ampliação do alcance das postagens. A Contratada deve desenvolver o serviço em diversas redes sociais, definidas junto à
				contratante, com a gestão exclusiva dos recursos.
			\end{enumerate}
	\end{enumerate}

7.3 Prospecção, planejamento, implementação, manutenção e monitoramento
de soluções de Comunicação Digital, no âmbito do contrato.

7.4 Criação e execução técnica de ações e/ou peças de Comunicação
Digital;

7.5 Criação, implementação e desenvolvimento de formas inovadoras de
Comunicação Digital, destinadas a expandir os efeitos de mensagens e
conteúdo do Conselho Regional de Psicologia São Paulo, em seus canais
proprietários e em outros ambientes, plataformas ou ferramentas
digitais, em consonância com novas tecnologias.

7.6 Diante desses fatores, a contratação de uma empresa de comunicação
digital é a solução mais eficiente e econômica para atender às demandas
crescentes do CRP~SP, garantindo a qualidade, eficácia e eficiência das
estratégias de comunicação adotadas pelo órgão.

\chapter{Justificativa para parcelamento}

8.1 Não haverá o parcelamento da solução em razão da economia de escala
a ser proporcionada pelo modelo adotado de contratação, e também pelo
ganho de efetividade na fiscalização. Ademais, o não parcelamento
permite a simplificação da gestão do contrato, que terá apenas um
fornecedor vencedor, o que garante também a padronização do serviço
ofertado.

\chapter{Demonstrativo dos resultados pretendidos}

9.1 A Comunicação digital atingiria os eixos dos Resultados Esperados da
Gestão do VXII Plenário do CRP~SP, especialmente no Eixo 2 (Processo de
Planejamento operacional do CRP 2024), que diz respeito ao diálogo da
psicologia com a sociedade brasileira e suas relações com a democracia e
direitos humanos.

9.2 Implementar soluções inovadoras, contribuindo assim para o alcance
dos objetivos de comunicação do CRP~SP, incluindo a promoção do diálogo
com a sociedade, o fortalecimento da marca e a disseminação de
informações relevantes para a categoria profissional.

9.3 Segmentação do público-alvo, garantindo um maior retorno sobre o
investimento e uma melhor mensuração dos resultados obtidos.

9.4 Promoção de ações que engajem a categoria na defesa dos Direitos
Humanos e do Estado Democrático de Direito, apoio na criação de
conteúdos que sensibilizem e mobilizem os psicólogos e psicólogas para a
importância da atuação profissional em consonância com os princípios dos
Direitos Humanos e da democracia, promovendo debates, eventos e
iniciativas  \textit{on-line} que estimulem a reflexão e a participação ativa da
categoria.

9.5 Contribuir para fortalecer o diálogo da psicologia com a sociedade
brasileira, ampliando o alcance das ações e promovendo o engajamento da
categoria na defesa dos Direitos Humanos e da democracia.

9.6 Disseminação de informações relevantes e atualizadas sobre ética na
Psicologia, favorecendo o aprimoramento contínuo dos profissionais e o
fortalecimento da categoria como um todo.

\chapter{Providências prévias ao contrato}

10.1 Não se vislumbra a necessidade de tomada de providências de
adequações para a solução ser contratada.

\chapter{Contratações correlatas/interdependentes}

11.1 Não se aplica

\chapter{Impacto ambiental}

12.1 A presente contratação não apresenta possibilidade de ocorrência de
impactos ambientais

\chapter{Viabilidade da contratação}

13.1 Esta equipe de planejamento declara viável esta contratação, diante
da fundamentação exposta neste documento, esta equipe de planejamento
declara viável esta contratação.

13.2 Este Estudo Técnico Preliminar evidencia que a contratação de
empresa especializada na prestação de Comunicação digital a fim de
atender às necessidades do CRP06 mostra-se tecnicamente viável e
necessária.

\chapter*{Equipe de planejamento da contratação}

\subsection*{Requisitante}
\rule[0.5ex]{\linewidth}{0.8pt}
\normalsize\textbf{Taís Aparecida de Souza}\\
\footnotesize Coordenadora de Comunicação\\
\rule[0.5ex]{\linewidth}{0.4pt}\normalsize
\vspace*{3\baselineskip}
	
\noindent\rule{28em}{0.4pt}\\
	\footnotesize São Paulo, 5 de abril de 2024\\\normalsize
	\rule{\linewidth}{0.8pt}
	\vspace{0.62\baselineskip}

\subsection*{Aprovação}

\rule[0.5ex]{\linewidth}{0.8pt}
\normalsize\textbf{Giane Del\textquotesingle Dono Rodrigues}\\
\footnotesize Gerente de Administração e Tecnologia da Informação\\
\rule[0.5ex]{\linewidth}{0.4pt}\normalsize
\vspace*{3\baselineskip}
	
\noindent\rule{28em}{0.4pt}\\
\\
	\noindent\rule{\linewidth}{0.8pt}

	\noindent\rule[0.5ex]{\linewidth}{0.8pt}
\normalsize\textbf{Talita Fabiano de Carvalho}\\
\footnotesize Conselheira presidente do CRP~SP\\
\rule[0.5ex]{\linewidth}{0.4pt}\normalsize
\vspace*{3\baselineskip}
	
\noindent\rule{28em}{0.4pt}\\
\\
	\rule{\linewidth}{0.8pt}

\end{document}
